\element{pdf}{
  \begin{questionmult}{pdf}\bareme{haut=2,p=0}
    Une particule de fluide
    % \begin{multicols}{2}
      \begin{reponses}
        \bonne{est un système mésoscopique}
        \bonne{contient en moyenne un un nombre de molécules constant}
        \bonne{suit le mouvement du fluide}
        \mauvaise{est un système microscopique}
      \end{reponses}
    % \end{multicols}
  \end{questionmult}
}

\element{deriveePart}{
  \begin{question}{deriveePart}\bareme{1}
    La dérivée particulaire s'écrit
    \begin{multicols}{4}
      \begin{reponses}
        \bonne{$\frac{\partial }{\partial t}+(\vv{v}.\vv{\text{grad}})$}
        \mauvaise{$\frac{\partial }{\partial t}+\text{div}\vv{v}$}
        \mauvaise{$\frac{\partial }{\partial t}-\text{div}\vv{v}$}
        \mauvaise{$\frac{\partial }{\partial t}+(\vv{\text{grad}})$}
      \end{reponses}
    \end{multicols}
  \end{question}
}

\element{constantes}{
  \begin{questionmult}{constantes}\bareme{haut=2,p=0}
    Quelles valeurs sont correctes (unités comprises)
    \begin{multicols}{3}
      \begin{reponses}
        \mauvaise{$\mu_{air}=\SI{1e3}{kg}$}
        \mauvaise{$\mu_{eau}=\SI{1e3}{g}$}
        \bonne{$\nu_{eau}=\SI{1e-6}{m^2.s^{-1}}$}
        \mauvaise{$\nu_{air}=\SI{1e-3}{Pl}$}
        \bonne{$\eta_{eau}=\SI{1e3}{Pa.s}$}
      \end{reponses}
    \end{multicols}
  \end{questionmult}
}

\element{debitMassique}{
  \begin{question}{debitMassique}
    Le débit massique se calcule avec
    \begin{multicols}{4}
      \begin{reponses}
        \bonne{$\iint_S\mu\vv{v}.\vv{dS}$}
        \mauvaise{$\iint_S\vv{v}.\vv{dS}$}
        \mauvaise{$\int_C \vv{v}.\vv{dS}$}
        \mauvaise{$\iiint_V\vv{v}.\vv{dS}$}
      \end{reponses}
    \end{multicols}
  \end{question}
}

\element{hydrostatique}{
  \begin{question}{hydrostatique}
    L'équation fondamentale de l'hydrostatique s'écrit (axe $z$ ascendant)
    \begin{multicols}{4}
      \begin{reponses}
        \bonne{$\vv{\text{grad}}P=\mu\vv{g}$}
        \mauvaise{$\vv{\text{grad}}P=-\mu\vv{g}$}
        \mauvaise{$\text{div}P=\mu\vv{g}$}
        \mauvaise{$\text{div}P=-\mu\vv{g}$}
      \end{reponses}
    \end{multicols}
  \end{question}
}

\element{tangentielle}{
  \begin{questionmult}{tangentielle}\bareme{haut=2,p=0}
    $\vv{\delta F}=\eta \frac{\partial \vv{v}}{\partial y}dS$
    % \begin{multicols}{5}
      \begin{reponses}
        \mauvaise{est la force de pression élémentaire}
        \mauvaise{est une force normale}
        \mauvaise{est la force exercée par une particule en $z$ sur une particule en $z+dz$}
        \mauvaise{est adaptée quand la vitesse est dirigée est selon $y$}
        \mauvaise{est une force conservative}
      \end{reponses}
    % \end{multicols}
  \end{questionmult}
}

\element{vitesseDebitante}{
  \begin{questionmult}{vitesseDebitante}\bareme{haut=2,p=0}
    La vitesse débitante
    % \begin{multicols}{1}
      \begin{reponses}
        \mauvaise{se mesure en \SI{}{m^3.s^{-1}}}
        \bonne{est $\frac{D_v}{S}$}
        \mauvaise{est le maximum de la vitesse sur une section}
      \end{reponses}
    % \end{multicols}
  \end{questionmult}
}

\element{reynolds}{
  \begin{questionmult}{reynolds}\bareme{haut=2,p=0}
    À bas Reynolds
    % \begin{multicols}{1}
      \begin{reponses}
        \bonne{la diffusion est prépondérante sur la convection}
        \bonne{le régime est laminaire}
        \mauvaise{les forces de frottements sont quadratiques}
        \bonne{$\Delta P=R_H D_v$}
        \mauvaise{le coefficient de trainée est indépendant du nombre de Reynolds}
      \end{reponses}
    % \end{multicols}
  \end{questionmult}
}